\SourceLeaf{389}%
nahme auf das Verhältniss 1 zu 40000, und für den Stich auf das Verhältniss 1 zu 80000, wonach die Anzahl der Blätter (von derselben Grösse wie oben) auf 208, die erforderliche Zeit auf 15 Jahr, die Kosten für die Arbeit, den Stich und den Abdruck von 3000 Exemplaren auf 4232000 Franken, endlich der Verkaufspreis jedes Blattes auf 7 Franken 50 Centimen, oder der ganzen Karte auf 1560 Franken veranschlagt werden. Die ganze Arbeit gehört zum Ressort des Depot, allein es ist dabei auf die Mitwirkung des Katasters gerechnet, obwohl aus dem Bericht nicht recht klar ist, in welchem Maasse: wie es scheint wird von dieser (von einem andern Ministerium abhängigen) Behörde das ganze Detail erwartet, so dass den Ingenieurs-Geographen bei der Aufnahme nur die trigonometrischen Arbeiten und die Höhenbestimmungen anheim fallen. Diese Abhängigkeit von einer andern Behörde, mit welcher kein recht harmonisches Zusammenwirken Statt zu finden scheint, (der nicht hinlänglichen Unterstützung von Seiten des Katasters wird das Fehlschlagen des ersten Plans beigemessen), könnte vielleicht dem gehofften raschen Fortgange dieser grossartigen Unternehmung sehr nachtheilig werden. Nach Vollendung der Arbeit soll noch ein grosses Repertorium geliefert werden, worin nicht bloss die numerischen Resultate für die Lage und Höhe der trigonometrischen Punkte, sondern auch alle Messungen auf welchen jene beruhen, bekannt gemacht werden sollen. Dadurch werden dann freilich alle Wünsche erfüllt werden. Allein so wie theils zu besorgen ist, dass dieser Zeitpunkt noch sehr weit entfernt sein möchte, theils auch in höhern wissenschaftlichen Beziehungen hauptsächlich nur die Dreiecke und Dreieckspunkte erster Ordnung das grösste Interesse darbieten, so können wir den lebhaften Wunsch nicht unterdrücken, dass man mit der vollständigen Bekanntmachung der Dreiecke erster Ordnung (welche bereits jetzt alle gemessen sind) nicht so lange zögern, sondern diese zum Besten der Wissenschaft sogleich liefern möchte. Die Freunde der höhern Geodaesie würden es um so dankbarer erkennen, wenn die künftigen Bände des Memorial diese Wünsche erfüllten. als sie, bei den bisher erschienenen Bänden, die in anderer Beziehung so gehaltreich sind, am wenigsten berücksichtigt worden sind, und in den wenigen theoretischen Dissertationen und Hülfstabellen keine Entschädigung für den Mangel an \textit{Thatsachen} finden, welche doch das Depot in so reichem Maasse zu geben im Stande wäre.
\begin{center}\rule{.3\textwidth}{.3pt}\end{center}
